
\documentclass{article}
\usepackage{colortbl}
\usepackage{makecell}
\usepackage{multirow}
\usepackage{supertabular}

\begin{document}

\newcounter{utterance}

\centering \large Interaction Transcript for game `cladder', experiment `full\_v1.5\_default', episode 9 with refineonly{-}Qwen3.5{-}27B.
\vspace{24pt}

{ \footnotesize  \setcounter{utterance}{1}
\setlength{\tabcolsep}{0pt}
\begin{supertabular}{c@{$\;$}|p{.15\linewidth}@{}p{.15\linewidth}p{.15\linewidth}p{.15\linewidth}p{.15\linewidth}p{.15\linewidth}}
    \# & \multicolumn{2}{c}{Player} && \multicolumn{2}{c}{Game Master} \\
    \hline

    \theutterance \stepcounter{utterance}  
    & & & \multicolumn{4}{p{0.6\linewidth}}{
        \cellcolor[rgb]{0.9,0.9,0.9}{
            \makecell[{{p{\linewidth}}}]{
                \texttt{\tiny{[P1$\langle$GM]}}
                \texttt{You are an expert in causal inference. The following question is not a typical commonsense query, but rather a meticulously designed question created by a professor specializing in causal inference, intended to assess the students' mastery of the course content.} \\
\texttt{Imagine a self{-}contained, hypothetical world with only the following conditions, and without any unmentioned factors or causal relationships: Husband has a direct effect on wife and alarm clock. Wife has a direct effect on alarm clock. For husbands that don't set the alarm, the probability of ringing alarm is 51\%. For husbands that set the alarm, the probability of ringing alarm is 49\%. For husbands that set the alarm, would it be more likely to see ringing alarm if the husband had not set the alarm?} \\
\texttt{Start your answer with 'Yes' or 'No', followed by additional reasoning or evidence to support your explanation.} \\
            }
        }
    }
     \\ \\

    \theutterance \stepcounter{utterance}  
    & \multicolumn{4}{p{0.6\linewidth}}{
        \cellcolor[rgb]{0.561,0.737,0.561}{
            \makecell[{{p{\linewidth}}}]{
                \texttt{\tiny{[P1$\rangle$GM]}}
                \texttt{No} \\
\\ 
\texttt{**Reasoning:**} \\
\\ 
\texttt{To answer this, we must distinguish between the **observed probability** (what we see in the data) and the **counterfactual probability** (what would have happened under a different condition).} \\
\\ 
\texttt{1.  **Analyze the Observed Data:**} \\
\texttt{    *   Let \$H\$ be the event the husband sets the alarm (\$H=1\$) or does not (\$H=0\$).} \\
\texttt{    *   Let \$A\$ be the event the alarm rings (\$A=1\$) or does not (\$A=0\$).} \\
\texttt{    *   The problem states:} \\
\texttt{        *   \$P(A=1 | H=0) = 0.51\$} \\
\texttt{        *   \$P(A=1 | H=1) = 0.49\$} \\
\texttt{    *   Observationally, the alarm is *less* likely to ring when the husband sets it.} \\
\\ 
\texttt{2.  **Analyze the Causal Structure:**} \\
\texttt{    *   The prompt specifies a self{-}contained world with only the following causal links: \$Husband \textbackslash{}rightarrow Wife\$, \$Husband \textbackslash{}rightarrow Alarm\$, and \$Wife \textbackslash{}rightarrow Alarm\$.} \\
\texttt{    *   There are no unmentioned factors or confounders.} \\
\texttt{    *   In such a world, the conditional probabilities \$P(A|H)\$ provided in the prompt represent the true causal effects (or the total effect of \$H\$ on \$A\$) because there is no back{-}door path to block. The probability of the alarm ringing given the husband sets it is exactly the probability of the alarm ringing under the intervention \$do(H=1)\$.} \\
\texttt{    *   Therefore, \$P(A=1 | do(H=1)) = 0.49\$ and \$P(A=1 | do(H=0)) = 0.51\$.} \\
\\ 
\texttt{3.  **Evaluating the Counterfactual:**} \\
\texttt{    *   The question asks: "For husbands that set the alarm (\$H=1\$), would it be more likely to see ringing alarm if the husband had *not* set the alarm (\$H=0\$)?"} \\
\texttt{    *   This asks us to compare the counterfactual outcome \$A\_\{H=0\}\$ for the subpopulation where \$H=1\$ against the actual outcome \$A\_\{H=1\}\$.} \\
\texttt{    *   Specifically, we are asking if \$P(A\_\{H=0\}=1 | H=1) > P(A\_\{H=1\}=1 | H=1)\$.} \\
\texttt{    *   Since the world is self{-}contained and the probabilities are defined by the causal structure, the probability of the alarm ringing if the husband had *not* set it is governed by the mechanism for \$H=0\$.} \\
\texttt{    *   The probability of the alarm ringing under the condition \$H=0\$ is 0.51.} \\
\texttt{    *   The probability} \\
            }
        }
    }
    & & \\ \\

    \theutterance \stepcounter{utterance}  
    & & & \multicolumn{2}{p{0.3\linewidth}}{
        \cellcolor[rgb]{0.9,0.9,0.9}{
            \makecell[{{p{\linewidth}}}]{
                \texttt{\tiny{[GM$|$GM]}}
                \texttt{No} \\
            }
        }
    }
    & & \\ \\

    \theutterance \stepcounter{utterance}  
    & & & \multicolumn{2}{p{0.3\linewidth}}{
        \cellcolor[rgb]{0.9,0.9,0.9}{
            \makecell[{{p{\linewidth}}}]{
                \texttt{\tiny{[GM$|$GM]}}
                \texttt{yes} \\
            }
        }
    }
    & & \\ \\

    \theutterance \stepcounter{utterance}  
    & & & \multicolumn{2}{p{0.3\linewidth}}{
        \cellcolor[rgb]{0.9,0.9,0.9}{
            \makecell[{{p{\linewidth}}}]{
                \texttt{\tiny{[GM$|$GM]}}
                \texttt{game\_result = LOSE} \\
            }
        }
    }
    & & \\ \\

\end{supertabular}
}

\end{document}
